% !PRINTED 341
\weberpage{341}
%<<CONT
holds, so long as $\xi, \eta$ belong to the domain (8). Then from (7) it follows that
\[
\beta' - \alpha' < (\beta - \alpha)\Theta.
\]
In the same way it follows, when we proceed in the same manner from the interval $(\alpha', \beta')$ to a narrower interval $(\alpha'', \beta'')$, that
\[
\beta'' - \alpha'' < (\beta' - \alpha')\Theta',
\]
in which $\Theta'$ has the same meaning for $\alpha', \beta'$ as $\Theta$ has for $\alpha, \beta$. But since $\alpha', \beta'$ belong to the domain (8), $\Theta'$ is not greater than $\Theta$, and $\Theta$ may also be put in place of $\Theta'$. It follows then:
\[
\beta'' - \alpha'' < (\beta - \alpha)\Theta^2,
\]
and so we conclude further
\[
\beta^{(\nu)} - \alpha^{(\nu)} < (\beta - \alpha)\Theta^\nu.
\]

The intervals therefore decrease at least as strongly as the terms of a decreasing geometric progression.

As an example may serve the equation:
\[
x^3 - 2x^2 - 2 = 0,
\]
which has a root between
\[
\alpha = 2{,}35 \quad \text{and} \quad \beta = 2{,}36
\]

From the formulae (4) one obtains for
\[
\beta' = 2{,}35931\ldots, \quad \alpha' = 2{,}359298\ldots,
\]
so that a value of the root accurate in the fourth decimal is
\[
2{,}3593
\]
For this value itself, as a more exact calculation shows, $f(x)$ is still negative, so that it can be taken for $\alpha'$. The next step of the approximation gives.
\[
2{,}359304.
\]

\begin{center}
\S.\,110.

\gesperrt{The approximation method of Daniel Bernoulli and related methods.}
\end{center}

The method for the approximate solution of an equation which is due to \gesperrt{Daniel Bernoulli} rests on the fact that, if one has a sequence of real quantities, the powers of the greatest among them will outweigh the equally high powers of the rest the more, the higher the powers are.
